% =========================================================
% Computational Linear Algebra --- Inverse Matrices
% (Strang \S{}2.5 / Johnston \S{}2.2)
% =========================================================

\section{Inverse Matrices}

% ---------------------------------------------------------
\subsection*{Definition and Properties}

\begin{propositionbox}{Definition (Invertible Matrix)}
A square matrix $A \in \mathbb{R}^{n \times n}$ is \emph{invertible}
(or \emph{nonsingular}) if there exists a matrix $A^{-1} \in \mathbb{R}^{n \times n}$
such that
\[
  AA^{-1} = A^{-1}A = I_n.
\]
The matrix $A^{-1}$ is the \emph{inverse} of $A$ and is unique.
A matrix with no inverse is called \emph{singular}.
\end{propositionbox}

\begin{theorembox}{Theorem (Properties of Inverses)}
When $A$ and $B$ are invertible $n \times n$ matrices:
\begin{enumerate}
  \item $(A^{-1})^{-1} = A$
  \item $(AB)^{-1} = B^{-1}A^{-1}$ \quad (order reverses)
  \item $(A^T)^{-1} = (A^{-1})^T$
  \item $(cA)^{-1} = \tfrac{1}{c} A^{-1}$ for $c \neq 0$
\end{enumerate}
\end{theorembox}

% ---------------------------------------------------------
\subsection*{\texorpdfstring{The $2 \times 2$ Case}{The 2 x 2 Case}}

\begin{tcolorbox}[colback=gray!6, colframe=gray!40, arc=2pt,
  left=6pt, right=6pt, top=4pt, bottom=4pt,
  title={\small\textbf{Formula ($2 \times 2$ Inverse)}},
  fonttitle=\small\bfseries]
\[
  \begin{pmatrix} a & b \\ c & d \end{pmatrix}^{-1}
  = \frac{1}{ad - bc} \begin{pmatrix} d & -b \\ -c & a \end{pmatrix},
  \qquad ad - bc \neq 0.
\]
The quantity $ad - bc$ is the determinant. The matrix is singular when $ad = bc$.
\end{tcolorbox}

% ---------------------------------------------------------
\subsection*{Change of Basis via Inverse}

\begin{remark*}[NURBS connection: converting between Bernstein and power basis]
In Piegl \& Tiller Exercise 1.14, the Bernstein basis functions of degree $n$
can be written as $[B_{0,n}(u), \ldots, B_{n,n}(u)] = [1, u, \ldots, u^n] M$,
where $M$ is an $(n+1) \times (n+1)$ matrix.
This means a B\'{e}zier curve has the matrix form
\[
  C(u) = [u^i]^T M [P_i],
\]
where $[u^i]^T$ is the monomial row vector, $[P_i]$ is the column of control points,
and $M$ encodes the change from power basis to Bernstein basis.

The inverse $M^{-1}$ converts the other direction: if you have a curve in
power basis form with coefficients $[\mathbf{a}_i]$, then the equivalent
B\'{e}zier control points are $[P_i] = M^{-1}[\mathbf{a}_i]$.

This is a direct application of the matrix inverse: $M$ converts one basis
to another, and $M^{-1}$ inverts the conversion.
\end{remark*}
